\scanpage{37}
\mspage{(2)}
s'interprète comme la classe de l'algèbre de quaternions définie par le
couple $(a,b)$. Appliquant un théorème dû à Witt $^{(4)}$, on a :

\textsc{Corollaire.} --- \textit{Le noyau de $\varphi$ est l'idéal engendré par les
éléments $e_a+e_b-e_{a'}-e_{b'}$ tels que $e_{a+b}=e_{a'+b'}$ et $a.b=a'.b'$.}

Soit $\mathrm H$ un anneau gradué à degrés positifs, et
$\widetilde{\mathrm H}=\mathbf Z\times(1+\widehat{\mathrm H}^{+})$
l'anneau filtré déduit de $\mathrm H$ par le procédé exposé par Grothendieck
dans $^{(2)}$. Alors la donnée d'un homomorphisme
\[
w^1:\quad\mathrm D(\mathrm A)\to\mathrm H^1
\]
tel que $w^1(e_a)w^1(e_b)=w^1(e_{a'})w^1(e_{b'})$ dès que
$(a,b,a',b')$, vérifient les conditions du corollaire, définit un homomorphisme
d'anneau filtré de $\mathrm Q(\mathrm A)$ dans $\widetilde{\mathrm H}$.

Prenons pour $\mathrm H$ l'algèbre de cohomologie $\mathrm H(\mathrm A)$ du groupe de Galois
de $\mathrm A$ à valeurs dans $\mathbf Z/2\mathbf Z$, la multiplication étant définie par le
cup-produit. Comme $\mathrm H^1(\mathrm A,\mathbf Z/2\mathbf Z)$ s'identifie à
$\mathrm D(\mathrm A)$ on a une application
\[
w:\quad\mathrm Q(\mathrm A)\to\widetilde{\mathrm H(\mathrm A)}.
\]
L'homomorphisme $w$ définit donc des applications
\[
w^i:\quad\mathrm Q(\mathrm A)\to\mathrm H^i(\mathrm A,\mathbf Z/2\mathbf Z).
\]
Soient $(\mathrm E,\mathrm Q)$ un module quadratique et $q$ sa classe dans
$\mathrm Q(\mathrm A)$; les $w^i(q)$ sont appelées les \textit{classes de
Stiefel-Whitney} de $(\mathrm E,\mathrm Q)$ (ou de $q$).
Comme $w$ est un homomorphisme d'anneaux, on a
\[
w^n(q+q')=\sum_{i+j=n}w^i(q)w^j(q').
\]
En particulier si $q$ est la classe du module
$\left[\mathrm A^p,\mathrm Q;\ \mathrm Q(x)=\sum_1^p a_jx_j^2\right]$, on a
\[
w^1(q)=\sum a_j;\qquad w^2(q)=\sum_{i<j}a_i.a_j.
\]
On interprète $w^1(q)$ comme le discriminant de $(\mathrm E,\mathrm Q)$,
et $w^2(q)$ de la façon suivante : si $rg(\mathrm E)=p$ est pair,
soit $\mathrm C(q)$ (resp. $\mathrm C_p$) l'algèbre de Clifford de
$(\mathrm E,\mathrm Q)$ [resp. de $p(\mathrm A,\mathrm I)$].
Il existe une extension galoisienne $\mathrm B$ de $\mathrm A$ telle que
$(\mathrm E,\mathrm Q)$ et $p(\mathrm A,\mathrm I)$ soient isomorphes sur $\mathrm B$.
Donc $\mathrm C(q)$ et $\mathrm C_p$ sont isomorphes sur $\mathrm B$.
On en déduit un élément $u(q)$ du groupe de Brauer de $\mathrm A$ par le procédé
habituel $^{(3)}$, et l'on a $u(q)=w^2(q)$. Si $p$ est impair, on a la même
interprétation en remplaçant $\mathrm C(q)$ par $\mathrm C^+(q)$
(resp. $\mathrm C_p$ par $\mathrm C_p^+$) $^{(1)}$.

Enfin on vérifie que $w$ est un homomorphisme du foncteur $\mathrm Q(\mathrm A)$
dans le foncteur $\mathrm H(\mathrm A)$. Ceci provient du fait que
$\mathrm H^1(\mathrm A,\mathbf Z/2\mathbf Z)$ est un foncteur covariant en $\mathrm A$.
\mspage{36}
